% EM structure and teaching revision 20260928
\subsection{端口等效的测定与适用边界}
线性有源二端网络可在端口写成 $U=U_{\mathrm{oc}}-IR_{\mathrm{th}}$，其中 $I$ 取网络向负载输出的方向。开路电压确定截距，短路电流在存在且有限时给出 $R_{\mathrm{th}}=U_{\mathrm{oc}}/I_{\mathrm{sc}}$。
只有独立源的电路，求等效电阻时才可把理想独立电压源短接、理想独立电流源断开；实际源内阻仍应保留。含受控源时必须保留控制关系，并在端口加测试电压或测试电流来求比值。
对 $R_{\mathrm{th}}>0$ 的被动负载 $R_L$，
\[
 P_L=\frac{U_{\mathrm{oc}}^2R_L}{(R_{\mathrm{th}}+R_L)^2},\qquad
 \frac{\dd P_L}{\dd R_L}=\frac{U_{\mathrm{oc}}^2(R_{\mathrm{th}}-R_L)}{(R_{\mathrm{th}}+R_L)^3}.
\]
最大功率对应阻值匹配，但此时等效内阻与负载各消耗一半功率。若追求输电效率，应使线路电阻远小于负载的等效阻值，目标与最大功率传输不同。
网络对称性也必须连同端口激励一起判断：图形上对称的两个节点，在不对称接入电源后未必等势。

